\section{A basis represented by a forest}\label{sec:forest}

The central combinatorial object is a pair $(\Fset,\Qset)$: an undirected
forest formed from selected precedence arcs, retaining each arc's original
orientation, and a set of root vertices. Every tree has one root except
for a single rootless tree. The theorem below establishes that this object
encodes exactly the nonsingular bases of the normalized LP. Its feasible
instances are then identified by a closure condition.

\begin{remark}[The forest is not new]
A forest carrying per-branch aggregates, with branches classified by the sign of
the aggregate they support, is the state maintained by the open-pit design
algorithm of \citet{LerchsGrossmann1965}: its \emph{normalized tree} hangs
branches from an added root, records the mass $b(T_q)$ of each branch, calls a
downward arc \emph{strong} when the mass it supports is positive, and pivots by
adding one arc and removing a branch's root arc. \citet{Hochbaum2001} sets that
algorithm out in this language and shows the tree determines a flow in linear
time. The same forest structure appears for the normalized ratio program in
\citet{Wuille2025}, and the revised-simplex algebra on rooted spanning forests
is developed in \citet{Yang2026} for a different linear program. What follows is
therefore a specialisation of the ordinary primal simplex to
\eqref{eq:standard}, not a new class of basis; the positioning document in
\path{literature/} states precisely what that leaves and what it does not.
\end{remark}

\subsection{Nonbasic slacks are the forest edges}
Let $\Nset$ denote the nonbasic variables in \eqref{eq:standard}, and define
\[
 \Fset=\{(i,j):s_{ij}\in\Nset\},\qquad
\Qset=\{i:y_i\in\Nset\}.
\]
The basic variables are recovered from the object by taking the complement:
\[
 \Bset=\{y_i:i\notin\Qset\}\ \cup\ \{s_e:e\notin\Fset\}.
\]
Thus a forest edge represents a \emph{nonbasic slack}, and a root represents
a \emph{nonbasic vertex variable}. Every remaining variable is basic, including
those whose value happens to be zero. The representation records which tight
constraints define the basis, even when additional constraints are also tight.

Setting the nonbasic variables to zero gives
\begin{equation}\label{eq:nonbasic}
 y_i=y_j\quad((i,j)\in\Fset),\qquad y_q=0\quad(q\in\Qset).
\end{equation}
Consequently, the vertex values are constant on each connected component of
$(\Vset,\Fset)$, with directions ignored only for connectivity.
The original orientation remains necessary when slack variables move.

\begin{theorem}[Forest characterization of bases]\label{thm:basis}
A set of $n-1$ nonbasic variables is the complement of a nonsingular basis
of \eqref{eq:standard} if and only if:
\begin{enumerate}[label=(\roman*)]
\item $\Fset$ is a spanning forest, with isolated vertices allowed;
\item each of its trees has at most one vertex in $\Qset$;
\item exactly one tree $T_*$ has no such vertex; every other tree has one.
\end{enumerate}
A vertex in $\Qset$ is called the root of its tree.
\end{theorem}
\begin{proof}
Eliminating all slacks leaves the normalization row and the $n-1$ equations
in \eqref{eq:nonbasic}, a square system on $y$. A cycle among forest edges
would make its edge-difference rows dependent. Likewise, two roots in one
component would be dependent modulo the path equations. Hence nonsingularity
requires an acyclic edge set and at most one root per component.
If the forest has $k$ trees, it has $n-k$ edges. Counting $n-1$ nonbasic
variables then requires $k-1$ roots, leaving exactly one tree rootless.
Conversely, the edge equations reduce $y$ to one scalar per tree. The
$k-1$ roots fix all but one scalar to zero, and normalization fixes the
remaining scalar uniquely because $w(T_*)>0$.
\end{proof}

At the resulting basic solution,
\begin{equation}\label{eq:basic-solution}
 y_i=\begin{cases}1/w(T_*),&i\in T_*,\\0,&i\notin T_*,\end{cases}
 \qquad s_{ij}=y_j-y_i,\qquad
 \rho=\vect{p}^T\vect{y}=\frac{p(T_*)}{w(T_*)}.
\end{equation}
We call $T_*$ the positive tree and all other trees zero trees.

\begin{corollary}[Feasibility and support]\label{cor:feasible}
The forest basis is primal feasible if and only if $T_*$ is a closure.
Every feasible basis has exactly one positive vertex level, and its support
is the connected closure $T_*$.
\end{corollary}
\begin{proof}
The only way a slack in \eqref{eq:basic-solution} can be negative is an arc
whose tail is in $T_*$ and whose head lies outside it. Excluding those arcs
is exactly the closure condition.
\end{proof}

This characterization concerns the normalized oracle, not every optimum of
the original capacitated LP. Moreover, a \emph{basic} variable can be zero.
An arc outside $\Fset$ has a basic slack, but that slack is zero whenever its
endpoints have the same value. These additional tight arcs account for many
degenerate pivots.

\subsection{A simple feasible initialization}
A nonempty DAG has a sink $v$, meaning that no arc leaves $v$ under our
prerequisite convention. Choose
\begin{equation}\label{eq:initial}
 \Fset=\emptyset,\qquad \Qset=\Vset\setminus\{v\},\qquad T_*=\{v\}.
\end{equation}
The basic variables are $y_v$ and all arc slacks. Set $y_v=1/w_v$ and all
other vertex variables to zero. The singleton $\{v\}$ is a closure, so this
basis is feasible. This construction works on disconnected graphs and needs
no artificial arcs or phase-I problem. For the numerical trace, choose the
smallest-index sink, vertex 1.

\subsection{Rooting and the orientation sign}
Root each zero tree at its root; root $T_*$ at any vertex $r_*$. For a tree
edge $e$, let $a$ be its parent endpoint and $b$ its child endpoint. Let
$D_e$ be the subtree rooted at $b$, including $b$, and put
\begin{equation}\label{eq:sign}
 \sigma_e=\begin{cases}
 +1,&e=(a,b)\text{ points from parent to child},\\
 -1,&e=(b,a)\text{ points from child to parent}.
 \end{cases}
\end{equation}
The slack equation implies
\begin{equation}\label{eq:parent-child}
 y_b-y_a=\sigma_e s_e.
\end{equation}
Thus a vertex value is its root value plus the signed sum of slacks on its
root-to-node path. On a general graph these signs depend on orientations,
not on path-length parity.
